\begin{center}\begin{tikzpicture}[font=\small,>=stealth,line join=round]
\draw[gray!40,rounded corners=3pt] (0,0) rectangle (14.2,10.2);\node at (7.1,9.6) {第一项高度更大：先增后减，只有最大值能取到};
\node[blue] at (7.1,8.65) {$\displaystyle f(x)=\sqrt{x^2+4}-\sqrt{(x-2)^2+1}$};\begin{scope}[shift={(3.1,4.6)},x=0.36cm,y=1cm]\draw[->,gray] (-6,0)--(14.5,0) node[right]{$x$};\draw[->,gray] (0,-2.8)--(0,3.4000000000000004) node[left]{$y$};\draw[gray,dashed] (-6,-2)--(14,-2);\node[left=3pt,fill=white,inner sep=1pt] at (0,-2) {$-2$};\draw[gray,dashed] (-6,2)--(14,2);\node[left=3pt,fill=white,inner sep=1pt] at (0,2) {$2$};\draw[blue,very thick,domain=-6:14,samples=240] plot(\x,{sqrt(\x*\x+4)-sqrt((\x-2)^2+1)});\draw[gray,dotted] (4,0)--(4,2.23606797749979);\fill[blue] (4,2.23606797749979) circle(2.3pt);\node[below=4pt] at (4,0){$4$};\node[blue,above=9pt] at (4,2.23606797749979) {$(4,\sqrt5)$};\end{scope}\node[anchor=west,align=left] at (9,4.5) {$a=2>b=1$\\[12pt]左端：趋近 $-2$\\[10pt]右端：趋近 $2$\\[12pt]最大值：$\sqrt5$\\[12pt]完整值域：$(-2,\sqrt5]$};
\node[text width=12.8cm,align=center] at (7.1,.65) {虚线是无限远处靠近的高度；不是有限定义域的端点。};
\end{tikzpicture}\end{center}
